\section{Smaller coefficients from projected crosses}
\label{sec:projected-crosses}
\label{sec:coefficients}

The coefficient compiler can discard information that the local computation
predicate never reads.  This gives a second, independent simplification of
the arithmetic construction.  Starting from the raw-machine blank tableau
of package~06~\cite{src06}, the full $3\times3$ window alphabet has
$1\,420\,903\,358\,357\,108\,792\,377$ members.  Five-cell crosses suffice,
and the previous-time cell need only retain its horizontal-boundary bit.
The resulting alphabet has exactly $87\,923\,161$ members.  The native
polynomial still uses four fixed integer coefficients, while the exact
binary length of $B-1$ falls from $5^{67}$ to $5^{28}$.

\subsection{The local language and its projection}

Use the same $439$ cell symbols as package~06.  Codes $0,1$ are headless
ordinary tape symbols.  Code $2(q+1)+s$ has tape symbol $s\in\{0,1\}$
and head state $q\in\{0,\ldots,215\}$.  The initialization symbols
$L,Q,R$ have codes $434,435,436$ and payloads
$(0,\varnothing),(0,0),(0,\varnothing)$, respectively.  The vertical
and horizontal boundaries $V,H$ have codes $437,438$.
The halt state $215$ is padded by stationary transitions for tableaux
only; the transition entering it retains its original write and move.

Let $\operatorname{Loc}$ be package~06's center predicate.  It checks
vertical-boundary propagation, separation of equal boundary types,
horizontal-boundary propagation away from vertical boundaries, the
initial phase path $L^uQR^v$, the terminal halt condition, nonescape of
the head, and the deterministic three-neighbor tape update.  In its
implementation \texttt{blank\_helical.local\_valid}, the four diagonal
corners are unused.  The previous-time neighbor is inspected only through
the test that it equals $H$; its other data are unused.

Set $a_0=439$.  A projected cross is
\[
 U=(p,\lambda,c,r,y)\in\{0,1\}\times\{0,\ldots,a_0-1\}^4.
\]
It is allowed when $\operatorname{Loc}$ accepts the $3\times3$ array with
center $c$, left $\lambda$, right $r$, following-time center $y$, zero
corners, and previous-time center $H$ if $p=1$ and $0$ if $p=0$.
This is an exact finite predicate.  Replacing the zero corners by any
symbols, or replacing the previous-time representative by any symbol
with the same $H$-indicator, leaves its truth value unchanged.
Write $\eta(c)=1$ if $c=H$ and $0$ otherwise.

For a cross $U$, its right neighbor $U'$, and its following-time
neighbor $U''$, impose just
\begin{equation}
 c=\lambda',\qquad r=c',\qquad
 \eta(c)=p'',\qquad y=c''.
 \label{eq:projected-overlap}
\end{equation}

\begin{lemma}[Exact projected-cross reconstruction]
For any cyclic length $N\ge1$ and any integer stride $h$, projected crosses
$U_i$ satisfying \eqref{eq:projected-overlap} with neighbors
$U_{i-1}$ and $U_{i-h}$ reconstruct a cell word $c_i$ whose local
computation predicate holds everywhere.  Conversely every valid cell
word gives such projected crosses.  This remains true when offsets
coincide, and requires neither $h\mid N$ nor coprimality.
\end{lemma}
\begin{proof}
The four equalities and their shifts give, with subscripts modulo $N$,
\[
 \lambda_i=c_{i+1},\quad r_i=c_{i-1},\quad
 y_i=c_{i-h},\quad p_i=\eta(c_{i+h}).
\]
Thus $U_i$ records exactly the data read by $\operatorname{Loc}$ at the
cell $c_i$.  Pull the word back to the plane by
$(x,t)\mapsto-x-ht\pmod N$.  The result is a valid plane with periods
$(N,0)$ and $(0,N)$.  The converse follows by taking the indicated data
from a cell word.  These equalities are literal even for coinciding
indices.
\end{proof}

The two distinguished tiles are
\[
 U_0=(1,434,435,436,1),\qquad
 U_1=(1,0,437,438,437).
\]
Both are allowed by the center predicate.  The first marks the initialized
head: K215's first transition writes $1$, moves right, and enters state
$93$, so the following-time center is the ordinary headless symbol $1$.
Every globally initialized $Q$ has this projected cross.  The second
tile cannot occur globally.  Its $p=1$ forces its predecessor to be $H$,
whereas its center is $V$.  Vertical-boundary propagation at that
predecessor would equate the vertical bits $0$ and $1$.

Order all allowed projected crosses as $U_0,U_1$, followed by the remaining
five-tuples in ordinary lexicographic order.  Package~06's helical padding
proof now applies through the reconstruction lemma: there is a marked
periodic word if and only if the raw K215 machine halts on blank input.
Completeness permits independently chosen sufficiently large widths and
heights, in particular powers of five.  Soundness does not require the
marked tile to be globally unique.  The canonical completeness word has
exactly one marked tile and no occurrence of the forbidden tile $U_1$.

\subsection{The exact tile count}

The inherited five-cell count is $38\,256\,607\,897$.  Its $13$ center
classes were rechecked against the literal local predicate by the
inherited finite census, considering every left/right pair for each
class.  Projection has constant fiber sizes.  An ordinary interior
center permits $438$ non-$H$ predecessors, all becoming one bit value.
An initialization center permits only $H$, so its count is unchanged.
At a boundary center the predecessor is unrestricted: $439$ choices
become two bit values.  The resulting exact counts are:
\begin{center}
\begin{tabular}{lr}
\toprule
Center class & Projected crosses\\
\midrule
$V$ & $384\,564$\\
$H$ & $767\,372$\\
Headless $0$ & $420\,300$\\
Halt with symbol $0$ & $223\,966$\\
Headless $1$ & $355\,508$\\
Halt with symbol $1$ & $223\,966$\\
Working head, writes $0$, moves left & $11\,570\,268$\\
Working head, writes $0$, moves right & $43\,205\,620$\\
Working head, writes $1$, moves left & $5\,046\,428$\\
Working head, writes $1$, moves right & $25\,720\,504$\\
Initial $L$ & $2\,118$\\
Initial $Q$ & $437$\\
Initial $R$ & $2\,110$\\
\midrule
Total $\kappa$ & $87\,923\,161$\\
\bottomrule
\end{tabular}
\end{center}
This count includes both distinguished tiles.  It performs no halting
test and does not discard locally allowed tiles on the basis of global
extendibility.

\subsection{Six native payload slots}

The native clause compiler copies the five coordinates of a tile and a
sixth, derived coordinate $\eta(c)$.  Its slots have the following exact
exponent and clause-index maps.  In every row $s$ ranges from zero to one
less than the slot size.
\begin{center}
\begin{tabular}{lrrl}
\toprule
Slot & Size & Radix exponent $e(t,s)$ & Clause exponent $j(t,s)$\\
\midrule
Previous $H$-bit & $2$ & $\kappa+10a_0+s$ & $1+s$\\
Left symbol & $a_0$ & $\kappa+5a_0+s$ & $3+s$\\
Center symbol & $a_0$ & $\kappa+4a_0+s$ & $3+a_0+s$\\
Right symbol & $a_0$ & $\kappa+3a_0+s$ & $3+2a_0+s$\\
Following symbol & $a_0$ & $\kappa+a_0+s$ & $3+3a_0+s$\\
Center $H$-bit & $2$ & $\kappa+7a_0+s$ & $3+4a_0+s$\\
\bottomrule
\end{tabular}
\end{center}
There are $4a_0+4$ payload positions.  The sixth coordinate is fixed by
the selected tile, so it is not an additional independent local choice.
Denote the six values of $U_i$ by $v_t(U_i)$, $0\le t\le5$.

Let $\ell\ge2$ be least with $2^\ell>\kappa+1$, and put $A_0=2^\ell$.
Define
\begin{align*}
 p_0&=\max\left(0,\left\lceil
       \frac{\kappa-2\ell-7a_0-10}{2}\right\rceil\right),&
 j_{\max}&=4a_0+8+p_0,\\
 \mu&=A_0-2+\sum_{j=1}^{j_{\max}}A_0^j,&
 m_0&=2\ell+11a_0+18+2p_0,\\
 n_{\rm pos}&=m_0+1,&
 n_{\rm dum}&=n_{\rm pos}-\kappa-4a_0-8,\\
 e_*&=\kappa+12a_0,&
 M_0&=\kappa+15a_0+n_{\rm dum}.
\end{align*}
Then $n_{\rm dum}\ge1$.  The allowed positions are the disjoint union
\[
 \mathcal E=\{0,\ldots,\kappa-1\}
 \cup\{e(t,s)\}
 \cup\{e_*,\ldots,e_*+n_{\rm dum}-1\}
 \cup\{M_0,3M_0,9M_0,27M_0\}.
\]
Their clause coefficients are
\begin{align*}
 c_i&=1+\sum_{t=0}^{5}A_0^{j(t,v_t(U_i))}
          +\sum_{u=1}^{4}A_0^{4a_0+4+u}
              &&(0\le i<\kappa),\\
 c_{e(t,s)}&=A_0^{j(t,s)},&
 c_{3^{u-1}M_0}&=A_0^{4a_0+4+u}\quad(1\le u\le4),
\end{align*}
and dummy positions have coefficient zero.  The bound
\[
 \sum_{e\in\mathcal E}c_e
 \le11(\kappa+4a_0+8)A_0^{4a_0+8}
\]
follows by counting the constant, six payload terms and four anchors of
each selector.

Set
\begin{align*}
 E_{\max}&=27M_0,&
 H_C&=E_{\max}+24M_0+3a_0+1,\\
 T_1&=H_C+2E_{\max}+a_0+1,&
 T_2&=T_1+4E_{\max}+2.
\end{align*}
Let $L_0$ be the least power of five at least $T_2+2E_{\max}+2$.
For a positive integer $x$, write $\operatorname{bits}(x)$ for its binary
length.  Put
\begin{align*}
 U_*&=4(n_{\rm pos}+2)\bigl(22(\kappa+4a_0+8)+6\bigr)+8,\\
 b_{\min}&=\max\{\ell(j_{\max}+1)+3,
            \ell(4a_0+8)+\operatorname{bits}(U_*)+1,4\}.
\end{align*}
Let $b$ be the least power of five at least $b_{\min}$, and set
\[
 \rho=2^b,\qquad d=bL_0,\qquad B=\rho^{L_0}=2^d.
\]

\subsection{The four fixed integers}

Define
\begin{align*}
 D_C={}&\rho^{3a_0}+\rho^{H_C+a_0}
       +\rho^{8M_0}+\rho^{24M_0}
       +\sum_{e\in\mathcal E}c_e
          (\rho^{T_1-e}+\rho^{T_2-e}),\\
 D_R={}&\rho^{H_C},\\
 F_0={}&\mu(\rho^{T_1}+\rho^{T_2})
       +\sum_{s=0}^{a_0-1}\rho^{\kappa+4a_0+s}
       +\sum_{s=0}^{1}\rho^{\kappa+10a_0+s}\\
     &+\sum_{u=4}^{5}\sum_{s=0}^{a_0-1}
             \rho^{H_C+\kappa+ua_0+s}
       +\rho^{9M_0}+\rho^{27M_0},\\
 M_C={}&B-1-\sum_{\substack{e\in\mathcal E\\ e\ne1}}\rho^e
              -2\rho^{e_*},\qquad
 F_{\rm nat}=F_0+4.
\end{align*}
The polynomial's four coefficient ports are exactly
\begin{equation}
 \boxed{\quad\beta=B-1,\qquad K=D_C+BD_R,\qquad
 M_C,\qquad M_F=F_{\rm nat}+B-1.\quad}
 \label{eq:projected-ports}
\end{equation}
The last coefficient is the shifted mask required by the arithmetic core.
All definitions are finite integer operations and bounded enumeration of
an explicitly given predicate; there is no free machine index or oracle.
These are four new integers and must be used together.

\subsection{Why the native compiler proof survives}

The four overlap tests in \eqref{eq:projected-overlap} are implemented by
the same two fixed translations as before.  Translation by $3a_0$ sends
the following-symbol slot to the center-symbol slot, and the center
$H$-bit slot to the previous $H$-bit slot.  Translation by $a_0$ sends
the right-symbol slot to the center-symbol slot and then to the left-symbol
slot.  The field masks test these four pairs; they contain $3a_0+2$
one-bit targets altogether.  Two further targets test the anchor pairs.

The two clause bands enforce zero or one selector at each cell, its six
matching one-hot payloads, and its four anchors.  Indeed, the low
$A_0$-digit test $A_0-2$, together with $A_0>\kappa+1$, permits only
selector totals zero and one.  Each higher one-bit clause requires its
payload bit to equal the corresponding selected-tile bit.  The added
center $H$-coordinate is thereby forced to be $\eta(c)$.
Parity of each sum of two Boolean payload bits is precisely equality,
which supplies \eqref{eq:projected-overlap}.  Occupancy propagates along
the horizontal neighbor relation from the marked origin to every cell.

There are no unintended convolution incidences.  The six payload intervals
are the explicit intervals in the preceding table; applying the two
translations gives exactly the four pairs just listed.  All other
nonanchor positions lie strictly below $M_0$.  At the anchor targets
$9M_0,27M_0$, the translations $8M_0,24M_0$ see exactly $M_0,3M_0$.
The minimum clause exponent $T_1-E_{\max}$ is larger than every copy or
anchor test exponent.  At a clause target, the only matching coefficient
comes from the identical source position.  The band separation prevents
cross-contamination.

More explicitly, the geometry satisfies
\[
 T_2-T_1>2E_{\max},\qquad
 L_0-(T_2-T_1)>E_{\max},\qquad T_2+E_{\max}<L_0.
\]
Thus an arbitrary rotation of native support, including the permitted
upper-dummy spill at $e_*+1$, cannot meet both clause bands.
Among allowed support positions the difference $18M_0$ occurs only for
the pair $(9M_0,27M_0)$.  There is no wrapped second possibility because
$L_0>45M_0$.  The inherited two-anchor argument therefore forces an
arbitrary binary rotation first to radix alignment and then to whole-cell
alignment.  The exceptional half-radix case has only one odd spill and
cannot satisfy the two anchor tests~\cite{masks}.

The radix bound gives
\[
 \rho\ge\max\left\{
  4(n_{\rm pos}+2)\left(2\sum_e c_e+6\right)+8,
  2\mu+4,16\right\}.
\]
Hence the inherited carry and representative bounds apply, including the
four-valued first dummy digit.  Also
\[
 \operatorname{popcount}(\mu)=\ell+j_{\max}-1,
 \qquad m_0=2\operatorname{popcount}(\mu)+3a_0+4.
\]
Consequently
\begin{gather*}
 0<M_C,F_{\rm nat}<B-1,\quad
 M_C\equiv2\pmod4,\quad F_{\rm nat}\equiv4\pmod8,\\
 \operatorname{popcount}(M_C)=d-m_0-1,\quad
 \operatorname{popcount}(F_{\rm nat})=m_0+1.
\end{gather*}
Their populations sum to $d$, exactly as required by the inverse-population
step of the arithmetic proof.  Selector $0$ is permitted at the marked
origin; selector $1$ is forbidden.  Its global impossibility was proved
above, so this exclusion removes no canonical accepting tableau.

Finally, $T_2-T_1\equiv2\pmod4$ and $\rho\equiv2\pmod5$ make every paired
clause contribution to $D_C$ cancel modulo five.  Because
$a_0\equiv3\pmod4$,
\[
 2D_C-D_R
 \equiv2\bigl(2+3\rho^{H_C}+1+1\bigr)-\rho^{H_C}
 \equiv3\pmod5.
\]
Thus the inherited dummy-bit subset argument still tunes the actual Pell
index to the intended time stride.  The required width, height and $d$
are powers of five.  The lower exponents $d$ change no logical step of
that argument.  The fixed low-field mask check also survives: the added
payload slots all have positive exponents, so the unit radix digit still
comes solely from the temporal neighbor's Start selector.  The canonical
word has exactly one such selector.

\subsection{Numerical dimensions and scope of verification}

The exact scalar parameters are
\begin{align*}
 \ell&=27,& p_0&=43\,960\,012,&j_{\max}&=43\,961\,776,\\
 m_0&=87\,924\,925,&n_{\rm dum}&=1,&e_*&=87\,928\,429,\\
 M_0&=87\,929\,747,&E_{\max}&=2\,374\,103\,169,\\
 H_C&=4\,484\,418\,415,&T_1&=9\,232\,625\,193,\\
 T_2&=18\,729\,037\,871,&b_{\min}&=1\,186\,967\,982.
\end{align*}
They give
\[
 b=5^{13},\quad L_0=5^{15},\quad d=5^{28}
       =37\,252\,902\,984\,619\,140\,625.
\]
The reduction in the binary length of $B-1$ is the exact factor $5^{39}$.
The other ports satisfy $M_C<B$, $M_F<2B$ and $K<B^2$; they too have
binary lengths bounded by $2d$.

The supplied implementation evaluates these scalar parameters, derives
the projected census from the checked $13$-class count, checks the
convolution incidences on small sparse fixtures, and proves the
actual-dimension support incidences by exact integer interval
intersections.  In particular the support audit does not sample the
$87\,923\,161$ selectors: it reasons about their entire interval.
It does not enumerate the full tile set, expand the four integers,
produce an accepting K215 computation, or materialize a full Pell witness.
The proof above is ordinary mathematics formalizable in ZFC.  The existing
K215 calibration then gives the requested equivalence; this section does
not add a kernel-checked Lean or Rocq theorem or claim a proof of
$\operatorname{Con}(\mathrm{ZFC})$.
